\section{Digital Citizen Reporting and e-Governance Systems}

Electronic governance (e-governance) is using the information and communication technologies to provide government services. This includes the support of government to citizens (G2C), government to business (G2B) , government to agencies (G2A) etc. By utilizing the e-governance platforms, everyone can gain access to whole range of public services. Additionally citizens can submit complaints and participate in governance decision making process. \cite{9587836}.

Digital citizen reporting system is one of the many use case of e-governance. "FixMyStreet" is a a such existing platform. It enable citizens to report issues for different issue categories. Sanitation, Traffic and Public Safety are some examples. Previous studies shows that these kind of systems improve governance administrative efficiency, civic participation and also they increase the transparency.\cite{oliveira2020smartcities}.

Even though there are many benefits of digital citizen reporting in local governance as described earlier, most of them still use centralized server based architecture. This centralized approach have its own limitations and  disadvantages. Susceptibility to data tempering, lack of transparency are the main problems that impact heavily on public trust. In addition to that, mandatory identity disclosure is discouraging public participation on this kind of reporting systems. Therefore these problems motivate the findings of new decentralized, transparent and immutable architectures that still preserve user privacy with accountability.




